\newcommand{\ConclusionContent}{%
{\RaggedLeft
\fontsize{21.5}{28}\selectfont
\textcolor{PosterPurple}{\textbf{دستاورد کمی:}}\ تنظیم نرخ پخت بر مبنای صدک \textbf{۷۱٪}، سالانه از هدررفت \textbf{۲۳۲ هزار پرس} غذا \mbox{(معادل \textbf{۸۷ میلیارد تومان})} با تقاضای برآورده‌نشده‌ی ناچیز (\textbf{\PosterNum{۰.۲۸٪}}) جلوگیری می‌کند.\par
\vspace{7pt}
\textcolor{PosterPurple}{\textbf{راهبرد مدیریتی:}}\ پایش رفتار مصرف و \textbf{شناسایی گروه‌های پراتلاف} جهت \textbf{تصمیم‌گیری‌های داده‌محور} به‌جای برش فله‌ای بودجه.\par
\vspace{7pt}
\textcolor{PosterPurple}{\textbf{محدودیت مدل:}}\ \textbf{بازه ۵ماهه داده‌ها} و \textbf{عدم پوشش نوسانات فصلی و تابستان}.\par
\vspace{7pt}
\textcolor{PosterPurple}{\textbf{افق و گام‌های آتی:}}\ \textbf{توسعه داشبورد تحلیلی}، \textbf{پیاده‌سازی آزمایشی (پایلوت)}، و \textbf{تعمیم مدل به سطح وزارتین} (عتف و بهداشت) جهت بهینه‌سازی هزینه‌های کلان سلف‌های سازمانی.\par
}
}
